% Generated draft: requires mathematical and source review.
\documentclass[11pt]{article}
\usepackage[margin=25mm]{geometry}
\usepackage{fontspec}
\setmainfont{Latin Modern Roman}
\usepackage{amsmath,amssymb,mathtools,unicode-math}
\setmathfont{Latin Modern Math}
\usepackage{enumitem,needspace,xurl,hyperref}
\setlength{\parindent}{0pt}
\setlength{\parskip}{5pt}
\setlength{\emergencystretch}{3em}
\newcommand{\N}{\mathbb{N}}
\newcommand{\Z}{\mathbb{Z}}
\newcommand{\Q}{\mathbb{Q}}
\newcommand{\R}{\mathbb{R}}
\newcommand{\C}{\mathbb{C}}
\newcommand{\F}{\mathbb{F}}
\newcommand{\A}{\mathbb{A}}
\newcommand{\PP}{\mathbb{P}}
\newcommand{\E}{\mathbb{E}}
\newcommand{\eps}{\varepsilon}
\newcommand{\dd}{\,\mathrm{d}}
\begin{document}
\section*{UnsolvedMath ID: 20003171}
\textbf{Source problem code:} catalogue entry
\subsection{Definitions and mathematical statement}
A Riemannian manifold is quasi-positively curved here if \(\sec\ge0\) everywhere and every tangent two-plane at at least one point has strictly positive sectional curvature. It is almost-positively curved in the sense relevant to the source if \(\sec\ge0\) everywhere and all two-planes are positive at every point of an open dense subset. The AIM problem asks for new examples of fundamental groups in these classes; ``positive on an open dense set'' is not interpreted as positivity of only one plane per point. Compactness and a reference set relative to which examples are new are not specified.

The supplied aid records Wilking's real family: for every \(n\ge3\), the real projectivized tangent bundle
\[
M_n=\PP_\R(T\R\PP^n)
\cong V_2(\R^{n+1})/(\Z/2)^2
\]
admits an almost-positive metric. Here \(V_2(\R^{n+1})\) is the space of ordered orthonormal two-frames, and the two sign changes act independently on the frame vectors. It is simply connected for \(n\ge3\), so the free quotient has
\[
\pi_1(M_n)\cong(\Z/2)^2.
\]
Let \(\lambda\) be the tautological line bundle over \(\PP_\R(T\R\PP^n)\). The aid reports
\[
w_1(TM_n)=n\,w_1(\lambda)\in H^1(M_n;\Z/2).
\]
Since \(\dim M_n=2n-1\), odd \(n\) gives nonorientable odd-dimensional examples. Synge's theorem then obstructs \emph{any} strictly positively curved metric on those members; it does not obstruct their reported almost-positive metrics. The restriction \(n\ge3\) is essential to the stated simply connected covering computation.

A second reported family comes from Kerr--Tapp: free cyclic quotients of the unit tangent bundle \(T^1S^6\) give complete quasi-positively curved eleven-manifolds with
\[
\pi_1\cong\Z/k\qquad\text{for every positive integer }k,
\]
including the trivial group at \(k=1\). The exact free actions and invariant metric construction are not supplied here; the fundamental-group statement uses the simply connected unit tangent bundle and the indicated free quotient.

The aid also records a general constraint: every \emph{closed} quasi-positively curved manifold has finite fundamental group. The Cheeger--Gromoll splitting argument would otherwise give a Euclidean factor in the universal cover and a zero-curvature mixed plane at every point, contradicting the strictly positive point. This statement retains closedness and is not silently extended to an arbitrary incomplete metric. The examples answer parts of the open-ended construction request, but do not classify all finite groups that occur or establish a new literature-novelty claim.
\subsection{Short English statement}
Find fundamental groups of quasi- or almost-positively curved manifolds. The supplied literature families realize \((\Z/2)^2\) and every finite cyclic group, while a closed quasi-positive manifold must have finite fundamental group; the real projective-tangent family also distinguishes almost-positive from strictly positive curvature.
\subsection{Sources}
\begin{itemize}
\item Problem 46, AIM workshop \emph{Manifolds with nonnegative sectional curvature}, \url{https://aimath.org/WWN/nnsectcurvature/nnsectcurvature.pdf}.
\item The research summary attributes the almost-positive projectivized-tangent family to Wilking and the unit-tangent-sphere cyclic quotients to Kerr--Tapp; it invokes Synge and Cheeger--Gromoll for obstructions. No full titles, dates, URLs, theorem numbers or free-action formulas are supplied.
\item Local UnsolvedMath registry, permanent ID 20003171, exact code \texttt{catalogue entry}; canonical location \texttt{aim-topology-notes.json}, notes[258], dated 14 August 2026; statement labeled exact.
\item Research attempt 1: partial result, problem open at the run, rigor label proved, novelty confidence low; registry status partially solved. The quotient and characteristic-class audits are machine-generated research aids, not peer reviewed; independent verification is required before citation.
\end{itemize}
\end{document}
